\documentclass[10pt,letterpaper]{article}
\usepackage{technical-user-manual}
\setDocTitle{Animation Timing: Sixteen-Week Curriculum}
\setDocSubtitle{Prerequisite-directed learning and practical evidence}
\setDocVersion{1.0}
\setDocStatus{Educational / Draft}
\setDocOwner{Game Development}
\setDocDate{2026-10-30}
\setDocSummary{Prerequisite-directed learning and practical evidence}
\begin{document}
\makedoctitle
\makedocsummary
\tableofcontents\clearpage
\section{Admission prerequisites}
Fundamentals: ratios, vectors, coordinate systems, basic plotting, and the ability to read a storyboard. Programming track: Python 3.11+ or C++20, a testing framework, and basic Git. Recommended supplementary mathematics: interpolation, derivatives, elementary kinematics, and numerical integration.
\section{Dependency graph}
\begin{verbatim}
Visual literacy + basic algebra
  -> Shot timing and FPS conversion
      -> X-sheets and spacing -> cycles and movement
      -> Storyboarding -> directing and editorial rhythm
  -> Vectors + derivatives -> kinematics -> integration
Python or C++ fundamentals -> animation data model -> unit tests
[spacing + kinematics + code tests] -> state machines / root motion
[storytelling + engine graph + testing] -> capstone / review
\end{verbatim}
\section{Sixteen-week sequence}
\begin{longtable}{p{.1\linewidth}p{.3\linewidth}p{.23\linewidth}p{.26\linewidth}}
\textbf{Week}&\textbf{Primary book reading}&\textbf{Dependency}&\textbf{Exit evidence}\\\hline
1 & Chapters 1--8 & storyboard literacy & six-shot animatic and direction notes\\
2 & Chapters 9--13 & ratios, frame units & frame conversions and 48-row X-sheet\\
3 & Chapters 14--17 & week 2 & distributed handoff and digital timeline\\
4 & Chapters 18--22 & vectors, forces & projectile contact-sheet experiment\\
5 & Chapters 23--26 & week 4 & articulated impulse chart\\
6 & Chapters 27--33 & weeks 2,4 & three spacing/hold alternatives\\
7 & Chapters 34--37 & week 6 & anticipation, overlap, spring settle\\
8 & Chapters 38--46 & weeks 6--7 & weight and repeat-action study\\
9 & Chapters 47--54 & weeks 6--8 & procedural effects and crowd variation\\
10 & Chapters 55--57 & cycles/contact & foot-contact walk with drift measurements\\
11 & Chapters 58--65 & week 10 & gait/flight comparison and aliasing review\\
12 & Chapters 66--72 & storyboards, week 2 & acting/lip-sync/beat-timing plan\\
13 & Chapter 73 & weeks 7,10 + code & interruptible animation state machine\\
14 & Chapters 74--77 & camera geometry & rig hierarchy and camera integration\\
15 & Chapters 78--82 & completed animatic & editorial revision and source synthesis\\
16 & Consolidation & all prior gates & playable capstone and formal scorecard\\
\end{longtable}
\section{Assessment gates}
Gate A: frame/time unit correctness. Gate B: expressive spacing and held actions. Gate C: numerical stability and reproducibility. Gate D: contact-safe gameplay animation transitions. Gate E: user experience evidence and engineering sign-off. Each gate requires an implementation or playblast, data sheet, automated test where applicable, reviewer comments, and revision record.
\section{Recommended use}
Do not reduce creative evaluation to automated metrics. Automated motion quality checks complement observed intent, readability and game-feel assessment.
\end{document}
